\unit{Question-type atlas, glossary and further practice}
\unitref{All tutorials, assignments and 2023--2025 examinations}
\subsection{Recognise a question before reaching for a formula}
The same few strategies recur throughout this course.
Use the table below to find the relevant teaching chapter
and then consult the longer worked-solutions booklets.
\begin{tabularx}{\linewidth}{@{}p{40mm}X@{}}
\toprule
\textbf{If the question asks about\ldots}&
\textbf{Start by doing this}\\\midrule
Coherent / thermal number counts&
Expand in \(\ket n\), form \(P_n\), evaluate
mean and factorial moments (Chapters 1--2)\\
HOM or photonic circuits&
Transform \emph{creation} operators and
normalise Fock states (Chapter 3)\\
Quadratures / squeezing&
State \([x,p]=2i\); compute means and
covariance, then apply the noise model (4--5)\\
Bell protocols and QKD&
Write the actual input state, gates,
measurement branches and sifted conditions (6--7)\\
Optical loss / postselection&
Include an environmental mode, trace or
project and normalise correctly (8--10)\\
Langevin cavity&
Commute Hamiltonian, Fourier transform and
apply the input--output boundary relation (9)\\
Atom--cavity coupling&
Choose conserved-excitation basis and
diagonalise a small block (11)\\
Quantum metrology&
Identify parameter generator and actual
measurement probabilities (12)\\
Hydrogen / nuclear size&
Separate \(R(r)Y_{lm}\); verify the radial
measure and use perturbation theory (13)\\
Configurations / LS terms&
Enforce antisymmetry; couple \(L,S,J\);
check the total state count (14)\\
3-j / transition amplitude&
Check triangle and \(m\) sums, then
Wigner--Eckart and parity (15,17)\\
Fine / hyperfine / Zeeman&
Choose \(L,S,J,I,F\), identify weak- or
strong-field basis, compute splitting (16)\\
Lifetime / multipoles&
Select allowed tensor rank and parity,
calculate dipole, sum all downward \(A\)'s (17)\\
Rabi pulse / laser cooling&
Distinguish coherent pulse area from
dissipative scattering force (18)\\
Atomic PV / EDM&
Find the symmetry-forbidden leading term,
then perturbatively mix states (19)\\
Cesium / optical clock&
Calculate hyperfine reference,
identify field-insensitive sublevels
and systematic shifts (20)\\
\bottomrule
\end{tabularx}

\subsection{Frequently confused concepts}
\textbf{Probability amplitude vs probability:}
a quantum state uses complex amplitudes;
only their modulus squared gives a probability.

\textbf{Coherent state vs pure state:}
a coherent state is one particular
pure oscillator state; many other
pure states are nonclassical.

\textbf{Mode overlap vs optical power:}
photon interference visibility depends
on mode indistinguishability, not just intensity.

\textbf{Relative phase vs global phase:}
a global factor \(e^{i\phi}\ket\psi\)
does not change measurement results,
but a relative phase between components does.

\textbf{Kraus channel vs conditional outcome:}
summed branches are trace preserving;
one selected detector outcome must
be renormalised and assigned a success probability.

\textbf{Angular momentum \(J\) vs term \(L,S\):}
total electronic \(J\) is the vector sum
of already coupled electronic orbital and spin moments.
Hyperfine \(F\) then adds nuclear \(I\).

\textbf{Allowed angular transition vs nonzero line strength:}
a 3-j symbol can pass selection rules
while the reduced atomic matrix element vanishes.

\textbf{Natural linewidth vs laser detuning:}
linewidth arises from decay/decoherence;
detuning is the frequency mismatch
between laser and atomic resonance.

\textbf{Clock accuracy vs stability:}
accuracy concerns systematic offset;
stability concerns fluctuations over averaging time.

\subsection{Physical constants and units that recur}
\[
\begin{array}{@{}ll@{}}
\hbar=h/(2\pi)&\text{reduced Planck constant}\\
a_0=5.29177\times10^{-11}\,\mathrm m&
\text{Bohr radius}\\
1E_h=27.2114\,\mathrm{eV}&\text{Hartree energy}\\
1E_h/(hc)=219474.63\,\mathrm{cm}^{-1}&
\text{atomic spectroscopy}\\
t_{\rm au}=2.418884\times10^{-17}\,\mathrm s&
\text{atomic unit of time}\\
\mu_B=e\hbar/(2m_e)&\text{Bohr magneton}\\
\alpha_{\rm fs}\simeq1/137.036&
\text{fine-structure constant}
\end{array}
\]
When a problem switches between wavelength,
wavenumber, frequency and angular frequency,
write each definition explicitly:
\[
\tilde\nu=1/\lambda,\quad
\nu=c/\lambda,\quad
\omega=2\pi\nu,\quad E=h\nu=\hbar\omega.
\]
The factors of \(2\pi\) are particularly
important in linewidth and Rabi calculations.

\subsection{Where to practise each technique}
The repository companion booklets are designed
to work together:
\begin{enumerate}
\item The \emph{15-page Course Overview} is the
fast exam-day map of key physical results.
\item The \emph{Worked Examples} booklet provides
twenty carefully selected full calculations.
\item The \emph{Assignment Solutions} booklet
develops derivations and numerical methods
associated with four assignment sheets.
\item The \emph{2023--2025 Past Exam Solutions}
booklet covers twelve full exam questions.
\end{enumerate}
In revision, read a mini-text chapter
and try the associated problem \emph{before}
looking at its worked solution.
List any remaining conceptual difficulty
in the repository's mistakes-and-gaps notes.
\begin{checkpoint}
An exam solution is strongest when it
connects a choice of model, a governing
operator, a calculation, a limiting-case
check and a physical interpretation.
\end{checkpoint}
